% !PRINTED 57
\weberpage{57}

%<<CONT
of analytic geometry) the \gesperrt{area of the period
parallelogram}.

By ranging congruent period parallelograms one against another,
the whole $u$-plane can be covered simply and without gaps.
Corresponding points of two different ones among these parallelograms
are the representatives of $u$-values that differ by integral
multiples of the periods, and that are called \gesperrt{congruent with respect
to the modulus $(\omega_1,\omega_2)$} or, when no doubt is possible,
\gesperrt{congruent} for short. The symbol of congruence is

\[ u \equiv u' \ (\mathrm{mod}\ \omega_1,\ \omega_2)\,, \]

and it says that $u'$ has the form $u+m_1\omega_1+m_2\omega_2$, if $m_1$, $m_2$
are integers.

% Source figure: Weber_Algebra_Bd3.pdf, PDF page 79, printed p.57, Fig.1; exact crop including caption and every original label.
% Former placeholder description retained: Four adjoining period parallelograms in the lattice of the periods $\omega_1,\omega_2$, with corner points $u_0$, $u_0+\omega_1$, $u_0+2\omega_1$, $u_0+\omega_2$, $u_0+\omega_1+\omega_2$, $u_0+2\omega_1+\omega_2$, $u_0+2\omega_2$, $u_0+\omega_1+2\omega_2$, $u_0+2\omega_1+2\omega_2$, and the interior points (joined by dashed lines) $u_1$, $u_2$, $u_3$, $u_4$.
\weberfigureimage{b3_i0078_fig1_source_20260831.png}{0.66}

Thus Fig.\,1 represents four period parallelograms laid one against
another; the points $u_1$, $u_2$, $u_3$, $u_4$ are congruent:

\[ u_2 = u_1+\omega_1\,,\quad u_3 = u_1+\omega_1+\omega_2\,,\quad u_4 = u_1+\omega_2\,, \]

and $u_1$, $u_2$, $u_3$, $u_4$ likewise form the corners of a period
parallelogram.

\gesperrt{The doubly periodic function $\varphi(u)$ has the same value at all
congruent points}, and the whole range of values of this function is
therefore exhausted \gesperrt{within a single period parallelogram}, provided
that of two opposite sides only \gesperrt{the one} is reckoned as belonging to
the parallelogram.

From our general assumptions we draw the
conclusion:

1. \gesperrt{There exists (apart from the constant) no doubly
periodic function that is free of points of discontinuity in the
period parallelogram;} for such a func-
%CONT>>
